\documentclass[12pt]{article}
\usepackage{latexsym, amscd, %graphicx, 
amssymb}
\newcommand{\SECT}[1]
		{\vspace{.3cm}
		$\!\!\!\!\!\!\!\!\!\!\!${\bf {#1}}
		\vspace{.2cm}}

\begin{document}


\begin{center}

{\large \bf Math 152, Fall, 2004, Workshop 6\par}
\vspace{.5cm}
{{\bf Honors Section} 
\par}
\end{center}

\begin{enumerate}

\item Four colonies of bacteria have the following 
population growth patterns: 

Colony II doubles its population every two days; 

Colony III triples its population every three days; 

Colony IV quadruples its population every four days; 

Colony V quintuples its population every ve days. 

\begin{enumerate}

\item Find the growth constant (the constant k in the equation 
$\frac{dP}{dt} 
= kP$) for each colony. 

\item Which colony is growing most rapidly? Which colony 
is growing most slowly? Which colonies, 
if any, are growing at the same rate? 

\end{enumerate}

\item Consider the differential equation $\frac{dy}{dx}
= \frac{x cos^{2}y}{2}$. Its direction field is plotted below. 

\begin{enumerate}

\item Using the direction field, sketch the solution curve that 
satisfies the initial condition $y(2) = 0$.

\item What is $\frac{dy}{dx}$
when $y = \frac{\pi}{2}$? How is this reflected in the direction 
field below? Sketch the 
solution of the differential equation satisfying the initial 
condition $y(0) = \frac{\pi}{2}$. 


\vspace{50em}



\item Find explicitly the general solution of the equation and 
then the solution $y = f(x)$ that satisfies 
the initial condition $y(2) = 0$. Using your calculator, plot this 
function. Include a separate sketch 
of this plot with your workshop (do not plot it on the direction field). 

\item Can the solution sketched in (b) be obtained from the general 
solution by choice of constant? 

\end{enumerate}

\item Find the general solution of the differential equation 
$e^{-y}\frac{dy}{dt}+\cos t=0$

\begin{enumerate}

\item Sketch the solutions corresponding to several values of the 
constant $C$ of integration (use the 
graphing function of your calculator). Does every value $C$ 
of the constant of integration correspond 
to a solution curve? If not, which values do? 

\item Do all the solutions have the same domain? Explain. 
Give an specific $C$ for which all real 
numbers t belong to the domain. 

\end{enumerate}

\end{enumerate}

\end{document}
